% =============================================================================
% Chapter 6 -- Geometry-Preserving Seed Selection
% =============================================================================

\chapter{Geometry-Preserving Seed Selection}
\label{ch:seed_selection}

This chapter presents the geometry-preserving seed selection strategy, the first original contribution of the circular oversampling component of this thesis.
The seed selection mechanism is a preprocessing step applied before any of the three proposed oversampling algorithms, ensuring that the minority instances used as generation seeds faithfully represent the distributional, geometric, and topological properties of the original minority class.
The Bhattacharyya Coefficient (BC)~\cite{bhattacharyya1943measure}, a well-established distributional similarity measure, serves as the primary marginal fidelity metric.


% ============================================================================
\section{Motivation: Why Seed Choice Matters}
\label{sec:seed:motivation}

In oversampling, the ``seed'' is the minority instance from which synthetic samples are generated.
In SMOTE and its variants, every minority instance is treated as an equally valid seed~\cite{chawla2002smote}.
However, research has repeatedly shown that not all minority instances are equally suitable:

\begin{itemize}[leftmargin=*]
  \item \textbf{Noisy seeds.} Minority instances in majority-dominated regions generate synthetic samples in the overlap zone, systematically worsening classifier performance~\cite{saez2015smote, batista2004study}.
  \item \textbf{Outlier seeds.} Minority instances far from the main cluster produce synthetic samples in empty feature-space regions, wasting the synthesis budget~\cite{han2005borderline}.
  \item \textbf{Redundant seeds.} Tightly clustered seeds produce near-duplicate synthetic samples, reducing distributional diversity~\cite{chawla2002smote, barua2014mwmote}.
  \item \textbf{Unrepresentative seed sets.} A seed set that fails to cover the full minority distribution produces augmented data whose marginal distributions deviate from the original, creating a train-test mismatch~\cite{he2009learning, fernandez2018learning}.
\end{itemize}

The goal of seed selection is to choose a \emph{subset} of minority instances that, when used as seeds, produces synthetic data maximally representative of the original minority class.
We formalise this through four complementary scoring metrics: the Bhattacharyya Coefficient (BC), JSD, AGTP, and a smoothness regulariser $Z$.

% ============================================================================
\section{Seed Selection Pipeline}
\label{sec:seed:pipeline}

The seed selection pipeline consists of four stages.

\subsection{Stage 1: Clustering the Minority Class}

The minority class is clustered into $K_{\mathrm{cluster}}$ groups using K-means~\cite{macqueen1967some, arthur2007kmeans}.
Clustering ensures candidate seed sets cover all minority sub-populations and provides the stratification structure for Stage~3.
This cluster-proportional stratification is the same principle used in K-Means SMOTE~\cite{douzas2018improving}, which showed that cluster-aware generation produces higher-quality synthetic distributions.

\subsection{Stage 2: PCA Projection}

The clustered minority class is projected to a $k_{pc}$-dimensional PCA subspace (default $k_{pc} = 2$)~\cite{jolliffe2002principal, bishop2006pattern}.
PCA removes correlation between features, making per-dimension histogram comparisons meaningful, and reduces the computational cost of scoring.

\subsection{Stage 3: Generating Candidate Seed Sets}

For each of $N_{\mathrm{candidates}}$ iterations, a candidate seed set $\mathcal{I}$ of size $m$ is generated by cluster-proportional stratified random sampling: from each cluster $c$, we draw $m_c = \lfloor (|C_c|/n) \cdot m \rceil$ seed indices without replacement.

\subsection{Stage 4: Score and Select}

Each candidate is evaluated using the composite criterion (Section~\ref{sec:seed:criterion}), and the candidate with the highest score is returned.

\begin{algorithm}[ht]
\caption{Geometry-Preserving Seed Selection}
\label{alg:seed_selection}
\KwInput{$\mathcal{S}_{\min}$: minority instances ($n \times d$); $m$: seeds to select; $N_c$: candidates; $K_{cl}$: clusters; $k_{pc}$: PCA components; $k_t$: topological $k$; $w_j$: JSD weight; $w_z$: smoothness weight}
\KwOutput{$\mathcal{I}^*$: selected seed indices; $s^*$: best score}

$\mathbf{X}_{pca} \leftarrow \text{PCA}(\mathcal{S}_{\min}, k_{pc})$\;
$\ell_1, \ldots, \ell_n \leftarrow \text{KMeans}(\mathbf{X}_{pca}, K_{cl})$\;
Compute proportional allocation: $m_c \leftarrow \lfloor (|C_c| / n) \cdot m \rceil$ for each cluster $c$\;
$s^* \leftarrow -\infty$; \quad $\mathcal{I}^* \leftarrow \emptyset$\;
\For{$t = 1$ \KwTo $N_c$}{
  \For{each cluster $c$}{
    $\mathcal{I}_c \leftarrow$ randomly sample $m_c$ indices from cluster $c$ (without replacement)\;
  }
  $\mathcal{I} \leftarrow \bigcup_c \mathcal{I}_c$; \quad $\mathbf{X}_s \leftarrow \mathbf{X}_{pca}[\mathcal{I}]$\;
  $s \leftarrow \BC(\mathbf{X}_{pca}, \mathbf{X}_s) + \mathrm{AGTP}(\mathbf{X}_{pca}, \mathbf{X}_s) - w_j \cdot \mathrm{JSD}(\mathbf{X}_{pca}, \mathbf{X}_s) - w_z \cdot Z(\mathbf{X}_s)$\;
  \If{$s > s^*$}{
    $s^* \leftarrow s$; \quad $\mathcal{I}^* \leftarrow \mathcal{I}$\;
  }
}
\Return{$\mathcal{I}^*$, $s^*$}
\end{algorithm}

% ============================================================================
\section{Bhattacharyya Coefficient for Marginal Similarity}
\label{sec:seed:bc}

The Bhattacharyya Coefficient~\cite{bhattacharyya1943measure} measures how faithfully a candidate seed set's marginal distributions reproduce those of the full minority class:
\begin{equation}
  \BC(\mathcal{X}, \mathcal{X}_s) = \frac{1}{k}\sum_{j=1}^{k} \sum_{b=1}^{B} \sqrt{p_b^{(j)} \cdot q_b^{(j)}},
  \label{eq:bc_seed}
\end{equation}
where $p_b^{(j)}$ and $q_b^{(j)}$ are normalised histogram proportions in bin $b$ of dimension $j$ for the full minority class and the seed set respectively, and $k$ is the number of features.
$\BC \in [0,1]$, with $\BC = 1$ indicating identical marginal distributions.
Higher BC implies greater overlap between the seed set and the original, ensuring that downstream oversampling generates from a representative distributional base.

% ============================================================================
\section{Jensen--Shannon Divergence Penalty}
\label{sec:seed:jsd}

JSD~\cite{lin1991divergence} serves as a penalty measuring distributional disagreement, complementing BC's agreement measure:
\begin{equation}
  \overline{\mathrm{JSD}}(\mathcal{X}, \mathcal{X}_s) = \frac{1}{k}\sum_{j=1}^k \mathrm{JSD}(P^{(j)}, Q^{(j)}).
  \label{eq:jsd_seed}
\end{equation}
Their combination prevents the degenerate case where matching broad support but different shapes yields a misleadingly high BC score.

% ============================================================================
\section{AGTP for Geometric-Topological Preservation}
\label{sec:seed:agtp}

AGTP evaluates the \emph{joint} structure of the two sets by combining geometric similarity (mean and covariance) and topological similarity ($k$-NN distance distributions).

\subsection{Geometric Similarity ($G_{\mathrm{sim}}$)}

\begin{equation}
  G_{\mathrm{sim}}(\mathcal{X}, \mathcal{X}_s) = \frac{1}{2}\left[\max\!\left(0, 1 - \frac{\|\bm{\mu}_{\mathcal{X}} - \bm{\mu}_{\mathcal{X}_s}\|}{\mathrm{Spread}(\mathcal{X}) + \varepsilon}\right) + \max\!\left(0, 1 - \frac{\|\mathbf{C}_{\mathcal{X}} - \mathbf{C}_{\mathcal{X}_s}\|_F}{\|\mathbf{C}_{\mathcal{X}}\|_F + \varepsilon}\right)\right],
  \label{eq:geom_sim}
\end{equation}
where $\bm{\mu}$ is the mean vector, $\mathbf{C}$ is the covariance matrix, and $\mathrm{Spread}(\mathcal{X}) = \frac{1}{|\mathcal{X}|}\sum_{\vx}\|\vx - \bm{\mu}_{\mathcal{X}}\|$.
This differs from FID~\cite{heusel2017gans} in that AGTP operates on raw features and uses the class spread as the denominator for numerical stability.

\subsection{Topological Similarity ($T_{\mathrm{sim}}$)}

\begin{equation}
  T_{\mathrm{sim}}(\mathcal{X}, \mathcal{X}_s) = \sum_{b=1}^B \min(p_b^{(\mathrm{kNN})}, q_b^{(\mathrm{kNN})}),
  \label{eq:topo_sim}
\end{equation}
where $p_b^{(\mathrm{kNN})}$ and $q_b^{(\mathrm{kNN})}$ are normalised histograms of $k$-NN distances within $\mathcal{X}$ and $\mathcal{X}_s$ respectively.
This measures whether the seed set has the same local spacing structure as the original minority class.

\subsection{Combined AGTP Score}

\begin{equation}
  \mathrm{AGTP}(\mathcal{X}, \mathcal{X}_s) = \frac{1}{2}\left[G_{\mathrm{sim}}(\mathcal{X}, \mathcal{X}_s) + T_{\mathrm{sim}}(\mathcal{X}, \mathcal{X}_s)\right] \in [0,1].
  \label{eq:agtp}
\end{equation}

AGTP is related in spirit to FID~\cite{heusel2017gans} and MMD~\cite{gretton2012kernel}, but differs by explicitly separating geometric and topological components, normalising each to $[0,1]$, and adding $k$-NN structure comparison that neither FID nor MMD provides.

% ============================================================================
\section{Smoothness Regulariser ($Z$)}
\label{sec:seed:z}

The smoothness regulariser penalises irregular seed spacing:
\begin{equation}
  Z(\mathcal{X}_s) = \frac{\sigma(\{d_1^{\mathrm{NN}}, \ldots, d_m^{\mathrm{NN}}\})}{\mu(\{d_1^{\mathrm{NN}}, \ldots, d_m^{\mathrm{NN}}\}) + \varepsilon},
  \label{eq:z_score}
\end{equation}
where $d_i^{\mathrm{NN}}$ is the distance from seed $i$ to its nearest seed neighbor.
Low $Z$ is desirable because evenly spaced seeds ensure full minority distribution coverage, while clustered seeds waste the synthesis budget on near-duplicate samples.

% ============================================================================
\section{Composite Selection Criterion}
\label{sec:seed:criterion}

The four metrics are combined into a single composite score:
\begin{equation}
  \mathrm{Score}(\mathcal{X}_s) = \underbrace{\BC(\mathcal{X}, \mathcal{X}_s) + \mathrm{AGTP}(\mathcal{X}, \mathcal{X}_s)}_{\text{representativeness}} - \underbrace{w_j \cdot \overline{\mathrm{JSD}}(\mathcal{X}, \mathcal{X}_s)}_{\text{divergence penalty}} - \underbrace{w_z \cdot Z(\mathcal{X}_s)}_{\text{spacing penalty}},
  \label{eq:seed_score}
\end{equation}
with default weights $w_j = 0.3$ and $w_z = 0.5$. The selected seed set is:
\begin{equation}
  \mathcal{X}_s^* = \argmax_{\mathcal{X}_s \in \mathcal{C}} \mathrm{Score}(\mathcal{X}_s).
\end{equation}

The criterion balances three objectives: (a)~maximize distributional representativeness; (b)~penalise distributional divergence; and (c)~penalise irregular spacing.

\begin{table}[ht]
\centering
\caption{Seed selection hyperparameters and default values.}
\label{tab:seed:hyperparams}
\begin{tabular}{llc}
\toprule
\textbf{Parameter} & \textbf{Description} & \textbf{Default} \\
\midrule
$N_{\mathrm{candidates}}$ & Number of candidate seed sets & 100 \\
$k_{pc}$ & PCA components for scoring & 2 \\
$K_{\mathrm{cluster}}$ & Clusters for stratified sampling & 5 \\
$k_t$ & $k$ for topological similarity & 5 \\
$w_j$ & JSD penalty weight & 0.3 \\
$w_z$ & Smoothness weight & 0.5 \\
$B$ & Histogram bins for BC/JSD/$T_{\mathrm{sim}}$ & 30 \\
\bottomrule
\end{tabular}
\end{table}

% ============================================================================
\section{PCA Projection Ablation Study}
\label{sec:seed:pca_ablation}

We ran seed selection with $k_{pc} = 2$ (default PCA) and $k_{pc} = d$ (no PCA: full feature space) on all 14 benchmark datasets.
Table~\ref{tab:seed:pca_ablation} reports the results.

\begin{table}[ht]
\centering
\caption{Effect of PCA projection on seed selection quality, averaged across 14 datasets (100 candidate evaluations per dataset, 5-fold CV with RF classifier). BC, AGTP, F1, AUC are higher-is-better; JSD, $Z$ are lower-is-better. Bold indicates the better value.}
\label{tab:seed:pca_ablation}
\begin{tabular}{lccccccc}
\toprule
\textbf{Setting} & \textbf{BC$\uparrow$} & \textbf{AGTP$\uparrow$} & \textbf{JSD$\downarrow$} & $Z\downarrow$ & \textbf{F1$\uparrow$} & \textbf{AUC$\uparrow$} \\
\midrule
No PCA ($k_{pc} = d$) & \textbf{0.8487} & \textbf{0.7772} & \textbf{0.0537} & \textbf{0.2959} & \textbf{0.8099} & \textbf{0.9474} \\
PCA ($k_{pc} = 2$)    & 0.8474 & 0.5657 & 0.0544 & 0.5224 & 0.7775 & 0.9433 \\
\bottomrule
\end{tabular}
\end{table}

BC is nearly identical with and without PCA (0.849 vs.\ 0.847) because the Bhattacharyya Coefficient is already robust to correlated features.
However, PCA substantially reduces AGTP ($-0.212$) and increases $Z$ ($+0.227$) because discarding all but two principal components destroys much of the joint geometric and topological structure.
As a result, downstream F1 is $0.032$ lower with PCA.

This finding motivates setting $k_{pc} = d$ (no PCA) as the default for AGTP and $Z$, or using $k_{pc} = \min(d, 5)$ to retain more structure without the full cost.
A detailed per-dataset breakdown and sensitivity analysis is provided in Chapter~\ref{ch:ablation}.

% ============================================================================
\section{Experimental Validation of Seed Selection}
\label{sec:seed:validation}

Table~\ref{tab:seed:quality_summary} reports the seed quality metrics for the selected seed sets compared to randomly generated seed sets, averaged across 14 benchmark datasets.

\begin{table}[ht]
\centering
\caption{Seed quality metrics: selected (geometry-preserving) vs.\ random (average of 100 random seed sets), averaged across 14 datasets.}
\label{tab:seed:quality_summary}
\begin{tabular}{lcccc}
\toprule
\textbf{Setting} & \textbf{BC$\uparrow$} & \textbf{AGTP$\uparrow$} & \textbf{JSD$\downarrow$} & \textbf{$Z\downarrow$} \\
\midrule
Selected (geometry-preserving) & 0.849 & 0.839 & 0.043 & 0.132 \\
Random (average) & 0.838 & 0.658 & 0.165 & 0.276 \\
\midrule
$\Delta$ & $+0.010$ & $+0.181$ & $-0.122$ & $-0.144$ \\
\bottomrule
\end{tabular}
\end{table}

The geometry-preserving selection achieves substantially higher BC and AGTP and lower JSD and $Z$ on every dataset.
This translates to classifier performance improvements reported in Chapter~\ref{ch:results}.
